% Canonical node 009 · A001; current section path 1.2.2.
% Canonical owner: A001
\cnnaNodeHeading{009}{A001}{Genesis seed $\star$}
\cnnaNodePlacement{1.2.2}{D0}

\paragraph*{Position and role.}
\texttt{A001} is an auxiliary root of the canonical DAG.  It is introduced at
this point because the exceptional constructor \texttt{C013} is the first node
that requires an explicit bootstrap argument.  A001 is not derived from
\texttt{C004A}; the two objects meet only as independent inputs of C013.

\paragraph*{Definition.}
The genesis-seed carrier is the singleton
\begin{equation}
  \mathcal S_\star:=\{\star\}.
  \label{eq:a001-singleton-seed}
\end{equation}
The token has no numerical, address, conductance, response, timing, geometric,
or dynamical component.  It is therefore neither a third free model parameter
nor an initial physical state.

\paragraph*{Formal encoding and countercheck.}
Python represents the carrier by an immutable zero-field dataclass.  Distinct
runtime instances compare equal because they have the same empty payload.  Lean
uses a one-constructor inductive type and proves
\begin{equation}
  \forall\eta:\mathcal S_\star,\qquad \eta=\star.
\end{equation}
The node-local test audits the empty field list and explicitly checks that the
canonical token exposes none of the downstream result fields.  A payload-bearing
seed would violate A001 even if it happened to be initialized with one value.

\paragraph*{Result and handoff.}
A001 closes only the existence of one explicit, information-free bootstrap
token.  Verified edge \texttt{E006} hands that token to \texttt{C013}.  The
claim that the first-birth result is independent of the supplied token belongs
to \texttt{T001}, not to A001.

\noindent\textbf{Primary code anchors.}
Python source and test: class \texttt{GenesisSeed}, lines 14--15, and
\texttt{TestGenesisSeed}, lines 13--24.  Lean: \texttt{GenesisSeed},
\texttt{canonical}, and \texttt{eqCanonical}, lines 12--27.  The Supplement
binds all five declaration-level anchors to their exact source hashes.
